\scanpage{72}
\mspage{II.20}
% archive sheet number: 71
% line 1; insertion at loop, written in box at foot
2.4. \inserted{faire d'abord une réduction au cas où $\pi=\pi_1(X,\unsure{\xi|\eta})$, en introduisant \deleted{\unsure{sur}} le rev. étale fini $X'$ de $X$ qui correspond à $\pi\subset\pi_1(X,\xi)$.}
% lines 1–6
On sait déjà que $V=\pi_1(\bar X)^{\mathrm{ab}}(\ell)$ est un module de t.f. sur $\unsure{\mathbb Z_\ell|Z_\ell}$, \deleted{ce qui} ce qui revient au $\unsure{=|\text{même}}$, que $H^1(\bar X,\mathbb Z/\ell\mathbb Z)$ est un groupe fini. \deleted{\unsure{(}} Ceci est vrai pour tt schéma de t.f. sur un corps \deleted{alg.} \inserted{sép.} clos, et se \inserted{peut} démontrer par réduction au cas des courbes, par le procédé du n° 3.
% lines 7–9
Lorsqu'on \deleted{\unsure{admet}} \inserted{sait} que $\pi_1(\bar X)^{\mathrm{ab}}(\ell)$ est de t.f. sur $\mathbb Z_\ell$, alors la relation à prouver, \deleted{\unsure{qui s'écrit}} s'écrit aussi \deleted{$\pi_1(\bar X)(\ell)_{\pi_0}$ fini i.e.}
% line 10
\[
\deleted{\unsure{(2.4)}}\qquad V_{\pi_0}\ \text{fini}\quad\text{i.e.}\quad (V\otimes_{\mathbb Z_\ell}\unsure{\mathbb Q_\ell|Q_\ell})_{\pi_0}=0,
\]
% line 11
et par dualité elle équivaut à la relation
% line 12
\[
(\check V)^{\pi_0}=0
\]
% lines 13–15
où $\check V=\operatorname{Hom}_{\mathbb Z_\ell}(\unsure{W|V|N},\mathbb Z_\ell)$. Or on a un isom canonique
\[
\check V\simeq\operatorname{Hom}(\pi_1(\bar X,\xi),\mathbb Z_\ell)\simeq H^1(\bar X,\mathbb Z_\ell),
\]
% lines 16–18
évidemment compatible avec l'action de $\pi_0$. Donc la conclusion de 1.1 s'écrit de façon équivalente, sous forme cohomologique
% line 19
\[
H_{\deleted{\unsure{1}}}^1(\bar X,\mathbb Z_\ell)^{\unsure{\bar\pi_0|\pi_0}}=0.\tag*{(2.4.1)}
\]
% lines lower left margin
\margin{faire une section 2.4\unsure{\ldots} spéciale pour cette \unsure{réduction}, \deleted{\unsure{pourrait déjà}} \uline{être \unsure{venu|revu}\unsure{e}}.}
